\section*{Hoofdstuk \ref{ch:b3:solid-state} --- Vastestofchemie: banden, defecten en halfgeleiders}

\begin{solution}{exo:b3:solid-state:1}
$(E - \alpha)/|\beta| = -2\cos(k\pi/7)$: $-1.802$, $-1.247$, $-0.445$, $+0.445$, $+1.247$, $+1.802$. Ze bestrijken
$3.604|\beta|$, al 90\,\% van de bandbreedte $4|\beta|$ van de oneindige keten.
\end{solution}

\begin{solution}{exo:b3:solid-state:2}
$2N$ ($N$ orbitalen, elk met twee elektronen). Natrium, met één elektron 3s
per atoom, vult zijn band s voor de helft: een metaal. Magnesium heeft er
twee en zou de band s precies vullen, maar de banden 3s en 3p overlappen,
zodat de hoogste bezette niveaus tot een gedeeltelijk gevulde
gecombineerde band behoren: ook een metaal.
\end{solution}

\begin{solution}{exo:b3:solid-state:3}
$\lambda \approx \qty{1239.8}{eV.nm}/E_g$: galliumfosfide \qty{549}{nm} (groen); galliumnitride \qty{365}{nm}
(nabij ultraviolet). Blauwe diodes gebruiken galliumnitride gelegeerd met
indium, waarvan de kloof kleiner is.
\end{solution}

\begin{solution}{exo:b3:solid-state:4}
$\ce{CaCl2} \xrightarrow{\ce{NaCl}} \mathrm{Ca_{Na}^{\bullet} + V_{Na}' + 2\,Cl_{Cl}^{\times}}$: twee kation- en twee anionplaatsen, zoals in NaCl; één Ca en
twee Cl; lading $+1 - 1 = 0$. $\ce{Y2O3} \xrightarrow{\ce{ZrO2}} \mathrm{2\,Y_{Zr}' + 3\,O_O^{\times} + V_O^{\bullet\bullet}}$: twee kation- en vier zuurstofplaatsen; twee Y
en drie O; lading $-2 + 2 = 0$.
\end{solution}

\begin{solution}{exo:b3:solid-state:5}
$T^{3/2} = 5196$ bij \qty{300}{K}, $kT = \qty{0.02585}{eV}$. Silicium: $\sqrt{N_cN_v} = \qty{2.42e19}{cm^{-3}}$, $\eu^{-1.125/0.05170} = \num{3.55e-10}$,
$n_i = \qty{8.6e9}{cm^{-3}}$. Germanium: $\sqrt{N_cN_v} = \qty{7.16e18}{cm^{-3}}$, $\eu^{-0.661/0.05170} = \num{2.80e-6}$, $n_i = \qty{2.0e13}{cm^{-3}}$. Verhouding
$\num{4.3e-4}$: de kleinere kloof van germanium geeft het ruim tweeduizend
keer meer \omterm{def:b3:solid-state:carriers}{ladingsdragers}.
\end{solution}

\begin{solution}{exo:b3:solid-state:6}
$n = N_D = \qty{1.0e16}{cm^{-3}}$; $p = n_i^2/n = \num{1.0e20}/\num{1.0e16} = \qty{1.0e4}{cm^{-3}}$.
\end{solution}

\begin{solution}{exo:b3:solid-state:7}
$E_F - E_i = kT\ln(n/n_i) = \qty{0.02585}{eV} \times \ln(10^6) = \qty{0.36}{eV}$.
\end{solution}

\begin{solution}{exo:b3:solid-state:8}
$kT = \qty{0.06894}{eV}$ bij \qty{800}{K}; $n/N = \eu^{-2.3/0.1379} = \num{5.7e-8}$; in een mol $\num{6.02e23} \times \num{5.7e-8} = \num{3.4e16}$
paren.
\end{solution}

\begin{solution}{exo:b3:solid-state:9}
De massa van een cel is $\rho a^3 = \qty{5.70}{g.cm^{-3}} \times (\qty{4.300e-8}{cm})^3 = \qty{4.532e-22}{g}$, d.w.z.\
$\qty{4.532e-22}{g} \times N_A = \qty{272.9}{g}$ per mol cellen, of \qty{68.23}{g} per zuurstofatoom. Dan
volgt uit $(1 - x) \times 55.8 + 16.0 = 68.23$ dat $1 - x = 0.936$: $x = 0.064$, $\mathrm{Fe_{0.936}O}$.
\end{solution}

\begin{solution}{exo:b3:solid-state:10}
$\ln(\sigma T)$: $-3.51$, $-1.30$, $0.36$ bij $1/T = 3.333$, $2.857$, $\qty{2.500e-3}{K^{-1}}$; helling tussen de eindpunten
$(0.36 + 3.51)/(-\qty{8.33e-4}{K^{-1}}) = -\qty{4.65e3}{K}$ (het middelste punt ligt op de rechte), dus
$E_a = \qty{4.65e3}{K} \times \qty{8.617e-5}{eV.K^{-1}} = \qty{0.40}{eV}$.
\end{solution}

\begin{solution}{exo:b3:solid-state:11}
$\tfrac12\ce{O2} \rightleftharpoons \mathrm{O_O^{\times} + V_M' + h^{\bullet}}$, $K = [\mathrm{V_M'}][\mathrm h^\bullet]/p(\ce{O2})^{1/2}$. Ladingsbalans $[\mathrm h^\bullet] = [\mathrm{V_M'}]$, dus
$[\mathrm h^\bullet]^2 = K\,p(\ce{O2})^{1/2}$ en $[\mathrm h^\bullet] \propto p(\ce{O2})^{1/4}$; de geleidbaarheid volgt de \omterm{def:b3:solid-state:carriers}{gaten}.
\end{solution}

\begin{solution}{exo:b3:solid-state:12}
$n/N = \sqrt{N_i/N}\,\eu^{-\Delta H_F/2kT} = \sqrt 2\,\eu^{-1.4/(2 \times 0.05170)} = 1.414 \times \num{1.32e-6} = \num{1.9e-6}$.
\end{solution}

\begin{solution}{pb:b3:solid-state:1}
\textbf{1.} $\ce{Y2O3} \xrightarrow{\ce{ZrO2}} \mathrm{2\,Y_{Zr}' + 3\,O_O^{\times} + V_O^{\bullet\bullet}}$.

\textbf{2.} Plaatsen: twee kationplaatsen en vier zuurstofplaatsen (drie
bezet, één leeg), de verhouding 1\,:\,2 van $\ce{ZrO2}$. Massa: 2 Y en 3 O aan elke kant. Lading: $2 \times (-1) + 2 = 0$.

\textbf{3.} Eén vacature per twee yttriumionen: een halve per yttriumion.

\textbf{4.} Per $(\ce{ZrO2})_{0.92}(\ce{Y2O3})_{0.08}$ zijn er 0.92 Zr en 0.16 Y, 1.08 kationen: $x = 0.16/1.08 = 0.148$,
$\mathrm{Zr_{0.852}Y_{0.148}O_{1.926}}$.

\textbf{5.} $(x/2)/2 = 0.0370$: 3.7\,\% van de zuurstofplaatsen is leeg.

\textbf{6.} $8 \times 0.0370 = 0.296$ vacatures per cel.

\textbf{7.} $a^3 = (\qty{5.14e-8}{cm})^3 = \qty{1.358e-22}{cm^3}$: $0.296/\qty{1.358e-22}{cm^3} = \qty{2.2e21}{cm^{-3}}$.

\textbf{8.} Elk oxide-ion moet over een barrière van \qty{1.0}{eV} naar een
naburige vacature springen; de fractie van de pogingen die slagen, $\eu^{-E_a/kT}$, groeit zeer snel met $T$.

\textbf{9.} Bij \qty{1073.15}{K} is $kT = \qty{0.09248}{eV}$: $A = \sigma T\eu^{E_a/kT} = 0.030 \times 1073.15 \times \eu^{10.81} = \qty{1.6e6}{S.K.cm^{-1}}$.

\textbf{10.} Bij \qty{973.15}{K}: $\sigma = (A/T)\eu^{-11.92} = \qty{0.011}{S.cm^{-1}}$.

\textbf{11.} Bij \qty{573.15}{K}: $\sigma = (A/T)\eu^{-20.25} = \qty{4.5e-6}{S.cm^{-1}}$.

\textbf{12.} $R = L/(\sigma S) = \qty{0.10}{cm}/(\sigma \times \qty{0.20}{cm^2})$: \qty{46}{\ohm} bij \qty{700}{\celsius}, $\qty{1.1e5}{\ohm}$ bij \qty{300}{\celsius}.

\textbf{13.} $R < \qty{1}{k\ohm}$ vereist $\sigma > \qty{5.0e-4}{S.cm^{-1}}$; oplossen van $(A/T)\eu^{-E_a/kT} = \num{5.0e-4}$ (door proberen)
geeft $T \approx \qty{760}{K}$, ongeveer \qty{490}{\celsius}.

\textbf{14.} Het uitlaatgas is koud na een start en bij lage belasting; het
verwarmingselement brengt de elektrolyt binnen enkele seconden op zijn
werktemperatuur, zodat de regeling werkt vanaf de start, wanneer de
meeste verontreinigende stoffen uitgestoten worden.

\textbf{15.} $\ce{O2 + 4e- -> 2O^2-}$; in de \omterm{def:b3:solid-state:kroger-vink}{notatie van Kröger-Vink} $\ce{O2} + \mathrm{2\,V_O^{\bullet\bullet} + 4\,e'} \longrightarrow \mathrm{2\,O_O^{\times}}$.

\textbf{16.} Zuurstof wordt aan de luchtkant gereduceerd en de oxide-ionen
worden aan de uitlaatkant weer tot zuurstof geoxideerd: de cel brengt $\ce{O2}$ over van de lucht naar het uitlaatgas, met
$\Delta_rG = RT\ln(p_{\mathrm{uitl}}/p_{\mathrm{lucht}})$ per mol $\ce{O2}$ en vier elektronen: $E = -\Delta_rG/4F = (RT/4F)\ln(p_{\mathrm{lucht}}/p_{\mathrm{uitl}})$.

\textbf{17.} $RT/4F = 8.3145 \times 973.15/(4 \times 96485) = \qty{0.02096}{V}$.

\textbf{18.} Nul: de twee drukken zijn gelijk.

\textbf{19.} $(RT/4F)\ln 10 = \qty{48.3}{mV}$ per decade.

\textbf{20.} $\qty{0.02096}{V} \times \ln(0.2095/0.010) = \qty{0.02096}{V} \times 3.04 = \qty{0.064}{V}$.

\textbf{21.} $\qty{0.02096}{V} \times \ln(0.2095/\num{1.0e-20}) = \qty{0.02096}{V} \times 44.49 = \qty{0.93}{V}$.

\textbf{22.} Rijk uitlaatgas bevat koolstofmonoxide, waterstof en
onverbrande koolwaterstoffen; op de platina-elektrode reageert de weinige
overgebleven zuurstof ermee,
en $p(\ce{O2})$ wordt bepaald door de evenwichten $\ce{2CO + O2 <=> 2CO2}$ en $\ce{2H2 + O2 <=> 2H2O}$, die bij \qty{700}{\celsius} een
uiterst kleine waarde overlaten.

\textbf{23.} $p = 0.2095\,\eu^{-0.45/0.02096} = \qty{1.0e-10}{bar}$.

\textbf{24.} Rond de stoichiometrische verhouding gaat het uitlaatgas over
van een kleine overmaat zuurstof naar een kleine overmaat CO en
waterstof: $p(\ce{O2})$ daalt met zo'n achttien
machten van tien voor een kleine verandering in brandstof, en de spanning
springt van minder dan
\qty{0.1}{V} naar ongeveer \qty{0.9}{V}. De motorsturing voegt brandstof toe als de spanning laag is en
neemt er weg als ze hoog is, en houdt zo het mengsel op het
stoichiometrische punt waar de driewegkatalysator CO, koolwaterstoffen en
stikstofoxiden tegelijk omzet.

\textbf{25.} In rijk uitlaatgas bij \qty{700}{\celsius} geeft de sensor ongeveer \qty{0.93}{V}.
\end{solution}
